% Research fragment for integration. No manuscript file is changed here.
% Kernel counterpart: PaperD.BiasedModes.diff_no_pattern and no_four_increments.

\begin{proposition}[Two modes for every initial bias]\label{prop:all-bias-modes}
Let $0<\varepsilon<1$, let $q\in[0,1]$ be the probability that the first
step is positive, and let $P_n(k)$ be the law of the number of positive
steps after $n\ge1$ steps. Extend $P_n$ by zero outside $\{0,\ldots,n\}$.
The first differences $d_n(k)=P_n(k+1)-P_n(k)$ satisfy
\[
 \operatorname{sp}\bigl(d_n(-1),d_n(0),\ldots,d_n(n)\bigr)
 \preceq (+,-,+,-).
\]
In particular, $P_n$ has at most two modal plateaus, including endpoint
modes. Moreover, there are no integers $i_0<\cdots<i_4$ for which
\[
 P_n(i_0)>P_n(i_1)<P_n(i_2)>P_n(i_3)<P_n(i_4).
\]
\end{proposition}

\begin{proof}
Put $a=1-2\varepsilon$ and $r_t(k)=\varepsilon+ak/t$. The forward equation
is
\[
 P_{t+1}(k)=(1-r_t(k))P_t(k)+r_t(k-1)P_t(k-1),
 \qquad P_1=(1-q,q).
\]
Subtracting adjacent equations gives
\begin{equation}\label{eq:bias-first-diff}
 d_{t+1}(k)=r_t(k-1)d_t(k-1)+(1-r_t(k+1))d_t(k),
 \qquad 1\le k\le t-1.
\end{equation}
Both coefficients are at least $\min(\varepsilon,1-\varepsilon)>0$:
$r_t(j)$ is a convex combination of $\varepsilon$ and $1-\varepsilon$
for $0\le j\le t$. We treat the two boundaries separately.

Suppose first that $q<1$. Since
$P_t(0)=(1-q)(1-\varepsilon)^{t-1}>0$, define
\[
 \xi_t=\frac{P_t(1)}{P_t(0)},\qquad
 b=\frac{1-2\varepsilon}{1-\varepsilon}<1,\qquad
 A=\frac{\varepsilon}{1-\varepsilon}>0.
\]
The boundary equation gives
\[
 \xi_{t+1}=(1-b/t)\xi_t+A,\qquad \xi_1=\frac q{1-q}.
\]
Thus $\delta_t=\xi_{t+1}-\xi_t$ obeys
\begin{equation}\label{eq:bias-edge-increments}
 \delta_{t+1}=\frac{t-b}{t+1}\delta_t+\frac A{t+1}.
\end{equation}
A nonnegative increment is followed by a positive increment. Consequently,
$\xi_t$ first decreases and then increases, allowing either stretch to be
empty. For $\ell_t=d_t(0)=P_t(0)(\xi_t-1)$ this proves
\[
 \operatorname{sp}(\ell_1,\ldots,\ell_N)\preceq(+,-,+).
\]
If $q\le1/2$, then $\xi_1\le1$ and the sharper bound is
$\operatorname{sp}(\ell_1,\ldots,\ell_N)\preceq(-,+)$.
In particular, a positive $\ell_i$ cannot be followed by a negative
$\ell_j$. Zeros do not affect either conclusion.

Write $r_t^*=d_t(t-1)$. Reflection of the walk gives
$r_t^*(q)=-\ell_t(1-q)$ and
$d_t^{,1-q}(t-1-k)=-d_t^{,q}(k)$.
For $0<q\le1/2$, the right boundary history therefore has pattern
contained in $(-,+,-)$. For $q=0$ it is nonpositive throughout: the
conditional count is at most $t-1$, so $r_t^*=-P_t(t-1)\le0$.

Fix $N$ and, for $1\le t\le N$, form the augmented row
\[
 Z_t=(\ell_N,\ldots,\ell_{t+1},
       d_t(0),\ldots,d_t(t-1),
       r_{t+1}^*,\ldots,r_N^*).
\]
For $q\le1/2$, its first row is contained in $(+,-,+,-)$.
Here is an explicit check that also handles zeros at
$\ell_1=r_1^*=2q-1$. If $Z_1$ contained four ordered signs $-,+,-,+$
and the second lay in the reversed left history, the first two would give
a positive left edge followed in time by a negative one, which is
impossible. Otherwise the last three lie in the right history and have
signs $+,-,+$. For $q>0$ their reflected negatives contradict the
$(-,+,-)$ exclusion for the left history at bias $1-q$; for $q=0$ they
contradict its nonpositive right history. This proves the asserted base
pattern, including $N=1$.

Every entry of $Z_{t+1}$ is a nonnegative combination of adjacent entries
of $Z_t$. The interior entries use \eqref{eq:bias-first-diff}; the two
new boundary entries and the remaining future edge values are copied
with coefficients $(1,0)$ or $(0,1)$. Such a step cannot introduce an
alternating strict-sign pattern: each nonzero entry has a parent of the
same sign, and parents of successive opposite signs are strictly ordered.
Thus $\operatorname{sp}(Z_N)\preceq(+,-,+,-)$. Reflection and sign
reversal preserve this pattern, proving the claim also for $q\ge1/2$,
including $q=1$.

Finally, extension by zero adds $d_N(-1)=P_N(0)\ge0$ on the left and
$d_N(N)=-P_N(N)\le0$ on the right; neither enlarges the pattern.
Every modal plateau contributes a positive-to-negative change, so there
are at most two. If the last displayed inequalities in the statement
held, each of their four intervals would contain an adjacent difference
of the same strict sign, by telescoping. These four differences would
have signs $-,+,-,+$, a contradiction.
\end{proof}

% The following is the short limit transfer. Its local-limit theorem is a
% published external input; the finite proposition above has no such input.
\begin{corollary}\label{cor:all-bias-limit-modes}
For $a\in(1/2,1)$ and $q\in[0,1]$, suppose the limiting density
$f_{q,a}$ has no constant interval. Then it has at most two local maxima.
\end{corollary}
\begin{proof}
The uniform local limit theorem of Trinas and Valle gives
\[
 \sup_{j\equiv n\;({\rm mod}\,2)}
 \left|\frac{n^a}{2}\Pr(S_n=j)-f_{q,a}(j/n^a)\right|\longrightarrow0.
\]
For fixed $x$, set $k_n(x)=\lfloor(n+n^a x)/2\rfloor$. Then
$0\le x-(2k_n(x)-n)/n^a<2n^{-a}$, and continuity gives
$n^aP_n(k_n(x))/2\to f_{q,a}(x)$.
If $x<y<z$ were local maxima, absence of a constant interval would give
$x<u<y<v<z$ with
$f_{q,a}(u)<\min\{f_{q,a}(x),f_{q,a}(y)\}$ and
$f_{q,a}(v)<\min\{f_{q,a}(y),f_{q,a}(z)\}$.
For example, choose a point below the smaller neighboring local maximum
in its one-sided maximizing neighborhood; the no-constant-interval
condition guarantees a strict decrease there.
The five corresponding lattice indices are strictly ordered for all
sufficiently large $n$, and convergence preserves these strict
inequalities. Proposition~\ref{prop:all-bias-modes} gives a contradiction.
\end{proof}
